% ================================================================
% SECTION 8: DATA MANAGEMENT
% ================================================================
\section{Data Management}
\label{sec:data_management}

\begin{center}
\code{bionetflux.core.domain\_data} ~~\textbar~~
\code{bionetflux.core.bulk\_data} ~~\textbar~~
\code{bionetflux.core.lean\_bulk\_data\_manager}
\end{center}

Three tightly related modules manage the per-element polynomial
coefficient arrays that represent the numerical solution.

% ------------------------------------------------------------------
\subsection{DomainData --- Lightweight Container}

\begin{center}
\code{bionetflux.core.domain\_data} \quad(52 lines)
\end{center}

A plain data-class that stores only the quantities needed for bulk operations on one domain:

\begin{lstlisting}[language=Python, caption={DomainData constructor}]
class DomainData:
    def __init__(self,
                 neq: int,
                 n_elements: int,
                 nodes: np.ndarray,
                 element_length: float,
                 mass_matrix: np.ndarray,
                 trace_matrix: np.ndarray,
                 initial_conditions: List[Optional[Callable]],
                 forcing_functions: List[Optional[Callable]]):
\end{lstlisting}

All array arguments are copied on construction.  Instances are created exclusively by the static factory in \code{BulkDataManager} (see below).

% ------------------------------------------------------------------
\subsection{BulkData --- Per-Domain Solution Storage}
\label{sec:bulk_data}

\begin{center}
\code{bionetflux.core.bulk\_data} \quad(449 lines)
\end{center}

\begin{structurebox}{BulkData internal layout}
Each \code{BulkData} instance stores an array
\[
  \mathtt{data} \in \mathbb{R}^{2 n_\mathrm{eq} \times n_\mathrm{el}}.
\]
For equation $e$ and element $k$ the two rows
$\mathtt{data}[2e:2e{+}2,\, k]$ hold the two bulk polynomial
coefficients $(\hat u_{0}, \hat u_{1})$ of that equation on that
element.  This representation follows the HDG convention of a P1
polynomial on the reference interval $[0,1]$.
\end{structurebox}

\subsubsection{Constructor}

\begin{lstlisting}[language=Python, caption={BulkData constructor}]
class BulkData:
    def __init__(self, problem: Problem, discretization: Discretization,
                 dual: bool = False):
\end{lstlisting}

\noindent On construction the class:
\begin{enumerate}
  \item Retrieves the trace matrix \code{T}, mass matrix \code{M}, quadrature matrix \code{QUAD} and quadrature nodes \code{qnodes} from an internally created \code{ElementaryMatrices(orthonormal\_basis=False)} instance.
  \item Allocates \code{self.data = np.zeros((2*neq, n\_elements))}.
\end{enumerate}

\subsubsection{set\_data --- Polymorphic Initializer}

\begin{lstlisting}[language=Python, caption={set\_data signature}]
def set_data(self, input_data, time: float = 0.0):
\end{lstlisting}

Dispatches to the \textbf{primal} or \textbf{dual} path depending on the \code{dual} flag.  Both paths accept three data formats:

\begin{longtable}{p{5.5cm}p{9cm}}
\toprule
\textbf{Input format} & \textbf{Action} \\
\midrule
\code{np.ndarray} of shape $(2 n_\mathrm{eq}, n_\mathrm{el})$ & Direct copy \\
\code{List} of $n_\mathrm{eq}$ callables $f(s,t)$ & \textbf{Primal:} evaluate at nodes, solve trace system.
  \textbf{Dual:} integrate via quadrature. \\
\code{np.ndarray} of size $n_\mathrm{eq}(n_\mathrm{el}{+}1)$ & Reconstruct coefficients from trace values \\
\bottomrule
\end{longtable}

\subsubsection{Query and Utility Methods}

\begin{longtable}{p{7cm}p{7.5cm}}
\toprule
\textbf{Method} & \textbf{Description} \\
\midrule
\code{get\_data() -> np.ndarray} & Return a copy of \code{self.data} \\
\code{get\_element\_data(k) -> np.ndarray} & Column $k$ of \code{data} (shape $(2 n_\mathrm{eq},)$) \\
\code{get\_trace\_values() -> np.ndarray} & Placeholder --- prints a warning \\
\code{compute\_mass(M) -> float} & $\sum_e \sum_k \mathbf{1}^T M \, \mathtt{data}[2e{:}2e{+}2, k]$ \\
\code{test() -> bool} & Built-in self-test (shape, NaN, bounds) \\
\bottomrule
\end{longtable}

% ------------------------------------------------------------------
\subsection{BulkDataManager --- Multi-Domain Coordinator}
\label{sec:bulk_data_manager}

\begin{center}
\code{bionetflux.core.lean\_bulk\_data\_manager} \quad(776 lines)
\end{center}

\begin{structurebox}{Design philosophy}
\code{BulkDataManager} stores only a list of \code{DomainData} objects.
Framework objects (\code{Problem}, \code{Discretization}, \code{StaticCondensation}) are passed as \emph{method arguments} rather than stored, keeping the manager lightweight and decoupled.
Every public method that accepts framework objects validates them against the stored domain data before proceeding.
\end{structurebox}

\subsubsection{Constructor and Factory}

\begin{lstlisting}[language=Python, caption={BulkDataManager initialization}]
class BulkDataManager:
    def __init__(self, domain_data_list: List[DomainData]):
        self.domain_data_list = domain_data_list

    @staticmethod
    def extract_domain_data_list(problems, discretizations,
                                 static_condensations) -> List[DomainData]:
        """Static factory: extract DomainData from framework objects."""
\end{lstlisting}

The factory method builds one \code{DomainData} per domain by:
\begin{enumerate}
  \item Calling \code{sc.build\_matrices()} to obtain mass and trace matrices.
  \item Extracting initial-condition and forcing callables from the \code{Problem} (\code{problem.u0[\ldots]} or \code{problem.force[\ldots]}).
\end{enumerate}

\subsubsection{Lifecycle Methods}

\begin{longtable}{p{8cm}p{7cm}}
\toprule
\textbf{Method} & \textbf{Description} \\
\midrule
\code{create\_bulk\_data(domain\_index, problem, discretization, dual=False) -> BulkData} & Create an empty \code{BulkData} for one domain \\
\code{initialize\_all\_bulk\_data(problems, discretizations, time=0.0) -> List[BulkData]} & Initialize all domains from initial conditions \\
\code{initialize\_bulk\_data\_from\_initial\_conditions(i, problem, disc, time) -> BulkData} & Initialize a single domain from ICs \\
\code{update\_bulk\_data(bulk\_data\_list, new\_data\_list)} & Overwrite bulk coefficients (validates shapes and NaN) \\
\bottomrule
\end{longtable}

\subsubsection{Computation Methods}

\begin{longtable}{p{8cm}p{7cm}}
\toprule
\textbf{Method} & \textbf{Description} \\
\midrule
\code{compute\_source\_terms(problems, discretizations, time) -> List[BulkData]} & Dual-mode integration of \code{problem.force} \\
\code{compute\_forcing\_terms(bulk\_data\_list, problems, discretizations, time, dt) -> List[np.ndarray]} & Combine mass-matrix $\times$ old solution + forcing $\times$ $\Delta t$ (implicit Euler) \\
\code{compute\_total\_mass(bulk\_data\_list) -> float} & Sum per-domain mass for conservation check \\
\code{compute\_mass\_conservation(bulk\_data\_list) -> float} & Alias for \code{compute\_total\_mass} \\
\bottomrule
\end{longtable}

\subsubsection{Utility Methods}

\begin{longtable}{p{7cm}p{7.5cm}}
\toprule
\textbf{Method} & \textbf{Description} \\
\midrule
\code{get\_num\_domains() -> int} & Number of domains \\
\code{get\_domain\_info(i) -> DomainData} & Inspect domain $i$ \\
\code{get\_bulk\_data\_arrays(bulk\_data\_list) -> List[np.ndarray]} & Extract raw arrays \\
\code{test(problems, discretizations, static\_condensations) -> bool} & Comprehensive self-test (9 sub-tests) \\
\bottomrule
\end{longtable}
